\documentclass[12pt,a4paper,oneside]{article}
\usepackage[brazil]{babel}
\usepackage[utf8]{inputenc}
\usepackage{olymp}

\setcounter{problem}{1}

\begin{document}

\begin{problem}{Notas (again)}{notas}{2}

Satisfeitos com seu programa da semana passada, que calculava a média das notas obtidas pelos alunos da turma, os professores de INF110 querem mais um auxílio para analisar o desempenho geral da turma: um programa que leia as notas da turma e informe a porcentagem de estudantes aprovados, reprovados e de exame final.

Estudantes com nota 60 ou mais são aprovados, com nota abaixo de 40 são reprovados, e os demais ($40 \leq \text{nota} < 60$) estão de exame final.

%\Notes
%
%Observações, se tiver.

\InputFile

A entrada contém duas linhas: a primeira contém somente um número inteiro $N$, que indica o total de alunos que fizeram a disciplina; a segunda contém $N$ valores reais, cada um deles indicando a nota final de um dos alunos. Restrições: todas as notas são valores reais entre 0 e 100 e há no máximo 1000 alunos.


\OutputFile

Seu programa deve gerar três linhas na saída: a primeira deve informar a porcentagem de alunos aprovados, a segunda de alunos reprovados e a terceira de alunos de exame final, seguindo o formato dado nos exemplos. Os valores devem ser escritos com uma casa decimal.

\Example

\begin{example}
\exmp{
5
92 55 62 84 70
}{
Aprovados: 80.0\%
Reprovados: 0.0\%
De exame final: 20.0\%
}%
\end{example}

\begin{example}
\exmp{
3
67.1 82.4 65
}{
Aprovados: 100.0\%
Reprovados: 0.0\%
De exame final: 0.0\%
}%
\end{example}


\begin{example}
\exmp{
6
39.9 40 60 100 59.9 0
}{
Aprovados: 33.3\%
Reprovados: 33.3\%
De exame final: 33.3\%
}%
\end{example}

%Comentários sobre o exemplo, se necessário.

\end{problem}

\end{document}
